\section{Conclusions}

The closed-FLRW gauge analysis resolves the principal ambiguity in the original bridge construction. Our \(n=1\) scalar harmonic maps to the exceptional \(N=2\) sector and is not retained as the ordinary propagating scalar carrier. The first physical scalar level is
\[
\boxed{n_\star=2}\qquad(\boxed{N_\star=3}).
\]
This first physical harmonic contains an exact sky-dipole subspace,
\[
T_{2,\rm dip}=A\sin[2\psi(z)](\hat{\mathbf n}\!\cdot\!\hat{\mathbf p}),
\]
when only the cross components \(A_{4i}\) are populated.

The corrected selective-response window is
\[
\boxed{\frac1{15}<\ell^2<\frac18},
\]
within which \(n=2\) is softened and every higher propagating mode \(n\ge3\) is stiffened. The next-level source-normalized response obeys
\[
\left|\frac{T_3/S_3}{T_2/S_2}\right|=\frac{8C_2}{15C_3}<\frac8{15}.
\]
The carried-forward supernova amplitude now calibrates the physical \(n=2\) dipole response. The R1 background and high-\(n\) growth predictions remain unchanged.

The revised signature is therefore
\[
\boxed{
\begin{gathered}
\text{first physical }S^3\text{ scalar mode }(n=2)\text{ enhanced}\\
\oplus\ \text{fixed-axis dipole subspace with }\sin2\psi\text{ envelope}\\
\oplus\ \text{suppressed }n\ge3\text{ susceptibility}\\
\oplus\ \text{near-GR high-}k\text{ growth}.
\end{gathered}}
\]
The next calculation is the full gauge-invariant \(n=2\) scalar--metric transfer matrix in the closed cubic-Galileon background.
